\documentclass[11pt,a4paper]{article}

% --- XeLaTeX: fuentes e idioma en Unicode ---
\usepackage{fontspec}
\usepackage[spanish]{babel}
\defaultfontfeatures{Ligatures=TeX}
% --- Matemáticas: primero lo clásico ---
\usepackage{amsmath,amssymb}
\usepackage{amsthm}
\theoremstyle{definition}
\newtheorem{definition}{Definición}[section]
\newtheorem{theorem}{Teorema}[section]
\newtheorem{remark}{Observación}[section]

% --- unicode-math SIEMPRE al final de los paquetes de math ---
\usepackage{unicode-math}
\usepackage{tikz}
\usetikzlibrary{arrows.meta,positioning}
\usepackage{graphicx}
\usepackage{booktabs}
\usepackage{siunitx}
\usepackage{natbib}
\usepackage{hyperref}
\usepackage{geometry}
\geometry{margin=2.5cm}

\title{Modelo matemático del sistema de bioconversión con \textit{Hermetia illucens}\\
\large Alimento -- Larvas -- Ambiente controlado: balances de masa, energía y gases}
\author{}
\date{}

\begin{document}

\maketitle
\tableofcontents
\newpage

\input{secciones/01_introduccion}
\section{Notación y convenciones}

% 2.1 Notación base de Eriksen (2022): A, B, L, rA, rB, rL,
%     rCO2m, rCO2B, rCO2L, a, m, YB, YL, Bmax, beta, instars
% 2.2 Notación adicional propia: N, S, T, w, c_CO2, c_CH4, c_O2, Q,
%     m_LC, R (agregadas), r_O2 (consumo de oxígeno, convención: positivo),
%     RQ, términos microbianos r^mic_CO2, r^mic_O2
% 2.3 Tabla de símbolos, unidades y dominios (escribir al final,
%     cuando todas las variables del modelo estén definidas)

\section{El sistema completo: población, sustrato y ambiente controlado}

\subsection{De una larva a una población}

El modelo de Eriksen describe la dinámica de los compartimentos de una larva
individual. Un sistema de bioconversión real contiene una población de
$N(t)$ larvas, cuyas tasas agregadas determinan los flujos de materia y
energía a escala del contenedor.

\begin{definition}[Tasa agregada de población]\label{def31}
Sea $r(t)$ una tasa por larva (asimilación, síntesis, producción de CO$_2$)
y $N(t)$ el número de larvas vivas en el contenedor. La \emph{tasa agregada}
de la población es
\begin{equation}
R(t) = N(t)\,r(t),
\label{eq:agregada}
\end{equation}
bajo el supuesto de población homogénea: todas las larvas comparten el mismo
estado $(A, B, L)$, es decir, la misma edad, peso y condición fisiológica.
\end{definition}

\begin{remark}\label{rem:homogeneidad}
El supuesto de población homogénea es razonable en crías sincronizadas
(larvas neonatas sembradas el mismo día, práctica estándar en producción),
pero se degrada conforme avanza el ciclo: la variabilidad individual en las
tasas de crecimiento genera una distribución de pesos que se ensancha con el
tiempo. Un refinamiento posible ---que dejamos como extensión--- es
reemplazar \eqref{eq:agregada} por la integral sobre la distribución de
pesos de la población, $R(t) = \int r(w)\,n(w,t)\,dw$, a costa de convertir
el modelo en una ecuación de estructura poblacional. En lo que sigue se
adopta \eqref{eq:agregada}, y la mortalidad se admite mediante $N(t)$
decreciente si los datos lo requieren.
\end{remark}

\subsection{El contenedor como volumen de control}

\begin{definition}[Sistema de bioconversión instrumentado]\label{def32}
El \emph{sistema completo} es el volumen de control delimitado por las
paredes del contenedor de cría, cuya frontera es atravesada por dos flujos
de aire ---\emph{inlet} y \emph{outlet}--- y en $t = 0$ por la carga inicial
del lote (larvas neonatas y sustrato). Su vector de estado se organiza en
tres grupos:
\begin{equation}
\mathbf{x} =
\underbrace{(A,\, B,\, L)}_{\text{larva promedio}}
\;\cup\;
\underbrace{(N,\, S)}_{\text{población y sustrato}}
\;\cup\;
\underbrace{(T,\, w,\, c_{\mathrm{CO_2}},\, c_{\mathrm{CH_4}},\, c_{\mathrm{O_2}})}_{\text{ambiente del aire}},
\label{eq:estado-completo}
\end{equation}
donde $S$ es la masa de sustrato disponible, $T$ la temperatura del aire,
$w$ la humedad absoluta y $c_i$ las concentraciones de los gases medidos.
\end{definition}

\begin{remark}\label{rem:frontera}
La elección de la frontera determina qué es flujo y qué es transformación
(cf.~Fig.~\ref{fig:balance}). Con la frontera en las paredes del contenedor,
la ventilación es el único flujo continuo que cruza la frontera
($J_{\mathrm{in}}$, $J_{\mathrm{out}}$), mientras que el consumo de
sustrato, el crecimiento larval, la respiración, la evaporación y la
metanogénesis son transformaciones internas $P$, $C$. El aire del
\emph{inlet} aporta O$_2$, vapor de agua y calor sensible a condiciones
controladas $(T_{\mathrm{in}}, w_{\mathrm{in}})$; el \emph{outlet} remueve
aire empobrecido en O$_2$ y enriquecido en CO$_2$, CH$_4$ y vapor, y es
precisamente en esa corriente de salida donde se instalan las mediciones
$\mathbf{y} = (T, H, \mathrm{CO_2}, \mathrm{CH_4}, \mathrm{O_2})$.
\end{remark}

\begin{remark}[Operación por lote]\label{rem:batch}
El sistema opera \emph{por lote}: en $t = 0$ se siembra una población de
$N_0$ larvas neonatas sobre una masa inicial $S_0$ de sustrato, se cierra el
contenedor y el ciclo transcurre sin recargas hasta que las larvas alcanzan
el estado de prepupa. No hay eventos discretos de \emph{carga}: el estado
$\mathbf{x}(t)$ evoluciona continuamente en todo el horizonte $[0, t_f]$.
La única conmutación programada es la de la ventilación: cada dos días el
\emph{inlet} se cierra durante un intervalo $[t_k^{c},\, t_k^{c}+\Delta]$
para estimar las tasas de producción de gases por acumulación (método de
cámara estática, Sección~6). Estas conmutaciones no alteran el estado, solo
activan o desactivan los términos de flujo de los balances
(cf.~Sección~4.2). Además de las mediciones de gas, el contenedor
está montado sobre una celda de carga que registra la masa total del lote,
$m_{\mathrm{LC}}(t) = S(t) + N(t)\bigl(A(t) + B(t) + L(t)\bigr) + m_f$,
donde $m_f$ es la tara (contenedor y estructura). La tasa de pérdida de masa
del lote obedece al flujo neto a través de la frontera:
\begin{equation}
-\dot{m}_{\mathrm{LC}} \;=\;
\dot{m}_{\mathrm{CO_2}} + \dot{m}_{\mathrm{CH_4}} + \dot{m}_{\mathrm{H_2O}}
\;-\; \dot{m}_{\mathrm{O_2,in}},
\label{eq:celda-carga}
\end{equation}
es decir, la masa de CO$_2$, CH$_4$ y vapor de agua que sale por el
\emph{outlet}, menos la masa de O$_2$ (y demás componentes) que el aire del
\emph{inlet} aporta y queda incorporada en esos productos. La
Ec.~\eqref{eq:celda-carga} es una restricción de consistencia global que
servirá de validación independiente de los balances de gases.
\end{remark}

\begin{remark}[Actividad microbiana del sustrato]\label{rem:microbiano}
El sustrato no es un almacén pasivo de alimento: alberga una microbiota que
respira ---consume O$_2$ y excreta CO$_2$---, genera calor y, en zonas
anaerobias del lecho húmedo, produce CH$_4$ por metanogénesis. En
consecuencia, las señales medidas en el \emph{outlet} son superposiciones de
la actividad larval y la microbiana,
\begin{equation}
r_{\mathrm{CO_2}}^{\mathrm{medido}} = N\,r_{\mathrm{CO_2}}
+ r_{\mathrm{CO_2}}^{\mathrm{mic}}(S, T, w),
\qquad
r_{\mathrm{O_2}}^{\mathrm{medido}} = N\,r_{\mathrm{O_2}}
+ r_{\mathrm{O_2}}^{\mathrm{mic}}(S, T, w),
\label{eq:co2-superpuesto}
\end{equation}
donde $r_{\mathrm{O_2}}^{\mathrm{medido}}$ denota la tasa total de consumo
de oxígeno (cf.~Ec.~\eqref{eq:rq}). El CH$_4$ proviene casi exclusivamente
del segundo proceso, lo que lo convierte en un indicador indirecto de la
anaerobiosis del sustrato. Nótese, además, que el cociente respiratorio
medido en el \emph{outlet} no es el cociente larval puro, sino
\begin{equation}
RQ_{\mathrm{medido}} \;=\;
\frac{N\,r_{\mathrm{CO_2}} + r_{\mathrm{CO_2}}^{\mathrm{mic}}}
     {N\,r_{\mathrm{O_2}} + r_{\mathrm{O_2}}^{\mathrm{mic}}},
\label{eq:rq-medido}
\end{equation}
de modo que la separación de ambas contribuciones ---un problema de
identificabilidad que se abordará en la sección de mediciones--- es
necesaria para interpretar $RQ$ en términos fisiológicos. Una vía
experimental directa es un lote testigo sin larvas, que mide
$r_{\mathrm{CO_2}}^{\mathrm{mic}}$ y $r_{\mathrm{O_2}}^{\mathrm{mic}}$ de
forma independiente.
\end{remark}

\begin{figure}[htbp]
\centering
\begin{tikzpicture}[
  font=\small,
  cont/.style={draw, very thick, rounded corners=6pt, minimum width=11.0cm,
               minimum height=4.8cm, fill=gray!5},
  flujo/.style={-{Stealth[length=3mm]}, very thick},
  comp/.style={draw, thick, rounded corners=3pt, fill=blue!8,
               minimum width=3.3cm, minimum height=1.55cm, align=center},
  sens/.style={draw, thick, fill=red!12, rounded corners=2pt,
               minimum width=1.0cm, minimum height=0.65cm, align=center},
  lbl/.style={fill=white, inner sep=1.6pt, font=\footnotesize, align=center},
  lblf/.style={above=2pt, midway, font=\footnotesize, align=center}
]

% Contenedor
\node[cont] (box) {};
\node[anchor=north west, xshift=8pt, yshift=-6pt] at (box.north west)
  {\textbf{Contenedor (volumen de control)}};

% Compartimentos internos (más arriba)
\node[comp] (sus) at ([xshift=-3.6cm,yshift=-0.4cm]box.center)
  {Sustrato\\[2pt] $S$};
\node[comp, fill=orange!12] (lar) at ([xshift=0cm,yshift=-0.4cm]box.center)
  {Población larval\\[2pt] $N$\\ $(A,B,L)$};
\node[comp, fill=green!8] (air) at ([xshift=3.6cm,yshift=-0.4cm]box.center)
  {Aire interno\\[2pt] $T,\, w$\\ $c_{\mathrm{CO_2}},\, c_{\mathrm{CH_4}},\, c_{\mathrm{O_2}}$};

% Intercambios internos: sustrato <-> larvas
\draw[{Stealth[length=2.2mm]}-{Stealth[length=2.2mm]}, thick, orange!85!black]
  (sus) -- node[lbl, above=6pt, midway]{consumo} (lar);

% Larvas -> aire: CO2, calor, vapor (flecha superior)
\draw[-{Stealth[length=2.2mm]}, thick, orange!85!black]
  (lar.15) -- node[lbl, above=6pt, midway]{CO$_2$, calor, vapor} (air.165);
% Aire -> larvas: O2 (flecha inferior)
\draw[-{Stealth[length=2.2mm]}, thick, orange!85!black]
  (air.195) -- node[lbl, below=6pt, midway]{O$_2$} (lar.345);

% Sustrato -> aire: evaporación, CH4, CO2 microbiano (curva inferior)
\draw[{Stealth[length=2.2mm]}-{Stealth[length=2.2mm]}, thick, orange!85!black, dashed]
  (sus.south) .. controls ([xshift=1.4cm,yshift=-1.1cm]sus.south)
    and ([xshift=-2.0cm,yshift=-1.2cm]air.south) ..
  node[lbl, pos=0.5]{evaporación, CH$_4$, CO$_2$ mic. / O$_2$ mic.} (air.south);

% Inlet (rótulo a la izquierda, dos líneas, flecha despejada)
\draw[flujo, green!50!black] ([yshift=0.9cm]box.west) ++(-2.0,0)
  -- node[pos=0, above=2pt, anchor=south west, font=\footnotesize, align=left]
     {inlet: $T_{\mathrm{in}},\; w_{\mathrm{in}}$,\\ $c_{\mathrm{O_2,in}},\; Q$}
  ([yshift=0.9cm]box.west);

% Outlet (rótulo a la derecha, flecha despejada)
\draw[flujo, red!70!black] ([yshift=1.3cm]box.east)
  -- node[pos=1, above=2pt, anchor=south east, font=\footnotesize, align=right]
     {outlet: $Q$}
  ++(2.0,0);

% Sensores en el outlet (debajo de la flecha)
\node[sens] (sensores) at ([xshift=1.0cm,yshift=0.55cm]box.east) {$\mathbf{y}$};
\draw[-{Stealth[length=2mm]}, thick, red!70!black, dashed]
  (sensores.north) -- ([xshift=1.0cm]box.east);
\node[below=3pt of sensores, align=center, font=\footnotesize]
  {$T,\ H,\ \mathrm{O_2}$\\ $\mathrm{CO_2},\ \mathrm{CH_4}$};

\end{tikzpicture}
\caption{El sistema completo de bioconversión instrumentado
(Definición~\ref{def32}). Dentro del volumen de control interactúan tres
grupos de compartimentos: sustrato, población larval y aire interno; las
flechas indican el sentido de cada intercambio (el O$_2$ fluye del aire
hacia larvas y microbiota; el CO$_2$, el calor, el vapor y el CH$_4$, hacia
el aire). La frontera es atravesada únicamente por los flujos de ventilación
(\emph{inlet/outlet}, caudal $Q$); las mediciones $\mathbf{y}$ se toman en
la corriente de salida y una celda de carga registra la masa total del
lote.}
\label{fig:sistema-completo}
\end{figure}

\subsection{Condiciones iniciales y horizonte del ciclo}

\begin{definition}[Ciclo de bioconversión]\label{def33}
Un \emph{ciclo de bioconversión} es la trayectoria del sistema
\eqref{eq:estado-completo} en el horizonte $[0, t_f]$, con condiciones
iniciales
\begin{equation}
N(0) = N_0, \quad S(0) = S_0, \quad
A(0) = A_0, \quad B(0) = B_0, \quad L(0) = L_0,
\label{eq:cond-iniciales}
\end{equation}
y ambiente en equilibrio con el \emph{inlet},
$T(0) = T_{\mathrm{in}}$, $w(0) = w_{\mathrm{in}}$,
$c_i(0) = c_{i,\mathrm{in}}$. El ciclo termina en $t_f$, definido como el
instante en que la fracción de prepupas en la población supera un umbral
$\pi_f$ convenido (típicamente $\pi_f \approx 0.5$), o equivalentemente
cuando la tasa agregada de asimilación cae por debajo del umbral de
mantenimiento de la población,
\begin{equation}
t_f \;=\; \inf\bigl\{ t > 0 \;:\; R_A(t) \;\leq\; N(t)\,m\,B(t) \bigr\}.
\label{eq:tf}
\end{equation}
En las condiciones de cría previstas, $t_f$ se espera dentro de la ventana
de $14$ a $16$ días. Cada ciclo se replica $n_{\mathrm{rep}} = 3$ veces en
condiciones nominales idénticas, lo que permitirá estimar la variabilidad
entre lotes de los parámetros del modelo (Sección~6).
\end{definition}

\begin{remark}[Calendario de cierres de ventilación]\label{rem:cierres}
El protocolo experimental incluye cierres del \emph{inlet} cada dos días,
en los instantes $t_k^{c} = 2k$ días, $k = 1, 2, \dots$, cada uno de
duración $\Delta$ por estimar (Sección~6). Durante un cierre, $Q = 0$ y las
concentraciones del aire interno evolucionan únicamente por la producción
interna, de modo que las pendientes $\Delta c_i / \Delta t$ miden
directamente las tasas $R_i$. El calendario $\{t_k^{c}\}$ y la duración
$\Delta$ son parte del protocolo: se conocen con exactitud y se registran
junto con las mediciones.
\end{remark}

\begin{remark}[Criterios de terminación y su medición]\label{rem:tf}
Los dos criterios de \eqref{eq:tf} no son equivalentes operativamente: la
fracción de prepupas requiere inspección visual o muestreo, mientras que el
criterio metabólico es \emph{observable en línea} ---con la salvedad de la
contribución microbiana al CO$_2$ (Observación~\ref{rem:microbiano}), que no
cesa al final del ciclo. Una firma de terminación más robusta es el cambio
de régimen del cociente CH$_4$/CO$_2$: al llegar la prepupa, las larvas
dejan de alimentarse y de remover el sustrato, cesa la bioturbación que
aireaba el lecho, se expanden las zonas anaerobias y el CH$_4$ aumenta
mientras el CO$_2$ larval cae a su componente de mantenimiento. Esta firma
es la que el gemelo digital puede evaluar en tiempo real sin intervenir el
lote. Para larvas neonatas de cría sincronizada, los valores iniciales
$A_0, B_0, L_0$ corresponden al peso seco promedio de una larva de un día,
repartido entre compartimentos según su composición bioquímica medida.
\end{remark}

\paragraph{Variables monitoreadas fuera del vector de estado.}
El mapa de observación $\mathbf{y}$ definido hasta aquí se restringe a las
variables que participan directamente en los balances del modelo:
$T$, $H$, CO$_2$, CH$_4$ y O$_2$. La instrumentación del sistema puede,
sin embargo, incluir sensores adicionales ---notablemente H$_2$S y NH$_3$---
cuyas señales no se incorporan como variables de estado. La razón es de
economía de modelado: incluir el H$_2$S exigiría abrir el balance de azufre
del sustrato, y el NH$_3$, el balance de nitrógeno con su equilibrio
amonio--amoníaco dependiente de pH; en ambos casos el costo en compartimentos
y parámetros adicionales es desproporcionado frente al valor que aportan al
objetivo central, que es describir la bioconversión de materia y energía.
Estos gases se clasifican entonces como \emph{variables de supervisión}:
mediciones en tiempo real que informan sobre la salud del proceso sin
alimentar las ecuaciones de estado.

\paragraph{Su rol en la capa de supervisión del gemelo digital.}
Aunque fuera del modelo dinámico, estas señales tienen una función bien
definida en el sistema monitoreado. El H$_2$S es un indicador de
putrefacción anaerobia de la fracción proteica del sustrato: su aparición
señala condiciones de falla ---exceso de humedad, compactación del lecho o
sustrato proteico sin consumir--- y opera como alarma temprana sobre el lote.
El NH$_3$, por su parte, es la vía gaseosa de pérdida de nitrógeno del
sistema: su concentración informa sobre la volatilización de amonio y,
por tanto, sobre la calidad final del frass como fertilizante. La distinción
introducida aquí ---variables de \emph{estado} que el modelo explica versus
variables de \emph{supervisión} que el sistema vigila--- es un principio
general del diseño: la frontera del modelo no la define lo que se puede
medir, sino lo que el modelo se propone explicar. Las extensiones a los
balances de nitrógeno y azufar quedan planteadas como trabajo futuro, una
vez validado el núcleo materia--energía del sistema.
\section{Supuestos del modelo}\label{sec:supuestos}

\subsection{El sistema en lazo cerrado}

Hasta aquí, las entradas $\mathbf{u}(t)$ se han tratado como forzamientos
exógenos. El sistema real, sin embargo, opera en \emph{lazo cerrado}: los
actuadores responden a las mediciones. Esta subsección formaliza esa
estructura antes de enunciar los supuestos del modelo.

\begin{definition}[Ley de control y consignas]\label{def41}
Sea $\mathbf{u}(t) = \bigl(Q(t), P_{\mathrm{cal}}(t), P_{\mathrm{hum}}(t)\bigr)$
el vector de entradas manipuladas ---caudal de ventilación, potencia de
calefacción y potencia de humidificación--- y
$\mathbf{r} = (T_{\mathrm{ref}}, H_{\mathrm{ref}}, \bar{c}_{\mathrm{CO_2}},
\bar{c}_{\mathrm{CH_4}})$ el vector de \emph{consignas} (valores de
referencia y límites de seguridad). Una \emph{ley de control} es un mapa
\begin{equation}
\mathbf{u}(t) = \pi\bigl(\mathbf{y}(t), \mathbf{x}_c(t); \mathbf{r}\bigr),
\label{eq:ley-control}
\end{equation}
donde $\mathbf{x}_c$ es el estado interno del controlador (por ejemplo, los
integradores de acciones PID). El \emph{sistema en lazo cerrado} es la
composición planta--controlador,
\begin{equation}
\dot{\mathbf{x}} = f\bigl(\mathbf{x}, \pi(g(\mathbf{x}), \mathbf{x}_c;
\mathbf{r}); \boldsymbol{\theta}\bigr),
\qquad
\dot{\mathbf{x}}_c = f_c\bigl(\mathbf{x}_c, g(\mathbf{x});
\mathbf{r}\bigr),
\label{eq:lazo-cerrado}
\end{equation}
con estado aumentado $(\mathbf{x}, \mathbf{x}_c)$ y exógenos reducidos a
las consignas $\mathbf{r}$ y las condiciones del \emph{inlet}.
\end{definition}

\begin{remark}[Estructura jerárquica del control]\label{rem:jerarquia}
El actuador $Q$ es compartido por todos los balances de gas: aumentar la
ventilación reduce CO$_2$ y CH$_4$ pero arrastra calor y humedad, en
conflicto con los lazos de $T$ y $H$. Se adopta por tanto una arquitectura
de dos niveles: lazos regulatorios de temperatura y humedad en el nivel
inferior, y una capa supervisora en el nivel superior que impone los
límites de seguridad de gases (incluidos H$_2$S y NH$_3$ como variables de
supervisión) y puede modificar las consignas o forzar ventilación máxima.
\end{remark}

\begin{remark}[Identificaci\'on en lazo cerrado]
\label{rem:id-lazo}
En~\eqref{eq:lazo-cerrado}, las entradas $u$ est\'an correlacionadas con las perturbaciones por construcci\'on: el controlador reacciona a ellas. Es la dificultad cl\'asica de la identificaci\'on en lazo cerrado~\citep{ljungSystemIdentification1999}: un estimador entrenado solo con datos de operaci\'on regulada puede aprender la ley $\pi$ en lugar de la fisiolog\'ia. En consecuencia, los datos de los cierres de ventilaci\'on (Observaci\'on~\ref{rem:cierres}) y de las ventanas con consignas en escal\'on desempe\~nar\'an el rol de datos de identificaci\'on, mientras que la operaci\'on regulada normal se reserva para validaci\'on (cf.~Secci\'on~6).
\end{remark}


\begin{definition}[Entrada de ventilaci\'on conmutada]
\label{def42}
Sea $\{t_k^c\}_{k \geq 1}$ el calendario de cierres programados (Observaci\'on~\ref{rem:cierres}) y $\Delta$ la duraci\'on de cada cierre. La ventilaci\'on opera en dos modos:
\begin{equation}
Q(t) =
\begin{cases}
Q_{\mathrm{op}}(t), & \text{modo ventilado: } t \notin \bigcup_k [t_k^c,\; t_k^c + \Delta],\\[4pt]
0, & \text{modo cerrado: } t \in [t_k^c,\; t_k^c + \Delta], \quad k = 1, 2, \ldots
\end{cases}
\label{eq:q-conmutada}
\end{equation}
donde en modo ventilado $Q_{\mathrm{op}}(t)$ est\'a determinado por la ley de control~\eqref{eq:ley-control}, y en modo cerrado la ventilaci\'on se anula por decisi\'on del protocolo, suspendi\'endose temporalmente la regulaci\'on.
\end{definition}

\begin{remark}[La conmutación no altera el estado]\label{rem:conmutacion}
A diferencia de un sistema híbrido con eventos de salto, en
\eqref{eq:q-conmutada} el estado $\mathbf{x}(t)$ es continuo en los
instantes de conmutación: lo que cambia es qué términos de los balances
están activos. En modo ventilado, los balances de gases y humedad incluyen
los flujos $J_{\mathrm{in}}$, $J_{\mathrm{out}}$ asociados a $Q$; en modo
cerrado esos términos se anulan y las concentraciones evolucionan solo por
producción interna,
\begin{equation}
\dot{c}_i = \frac{R_i(t)}{V_{\mathrm{aire}}},
\qquad t \in [t_k^{c},\, t_k^{c} + \Delta],
\label{eq:modo-cerrado}
\end{equation}
con $V_{\mathrm{aire}}$ el volumen de aire del contenedor. La
Ec.~\eqref{eq:modo-cerrado} es la base del método de cámara estática~\citep{hutchinsonImprovedSoilCover1981,ermolaevGreenhouseGasEmissions2019}: las
pendientes de acumulación miden directamente las tasas $R_i$
(Sección~\ref{sec:mediciones}).
\end{remark}

\begin{remark}[Duración del cierre]\label{rem:delta}
La duración $\Delta$ resuelve un compromiso entre detectabilidad y
seguridad. Por abajo, la acumulación debe superar con holgura la resolución
$\sigma_i$ de cada sensor:
\begin{equation}
\Delta \;\geq\; \max_i \frac{10\,\sigma_i\, V_{\mathrm{aire}}}{R_i^{\min}},
\label{eq:delta-min}
\end{equation}
evaluada en el período de menor actividad metabólica (inicio del ciclo).
Por arriba, el O$_2$ no debe caer por debajo de un piso seguro
$\underline{c}_{\mathrm{O_2}}$:
\begin{equation}
\Delta \;\leq\; \frac{\bigl(c_{\mathrm{O_2}}(t_k^{c})
- \underline{c}_{\mathrm{O_2}}\bigr)\, V_{\mathrm{aire}}}{R_{\mathrm{O_2}}^{\max}}.
\label{eq:delta-max}
\end{equation}
Ambas cotas son calculables \emph{antes} del experimento a partir del
modelo, y ajustables día a día por el gemelo digital conforme $R_i$
evoluciona (Sección~\ref{sec:estimacion}).
\end{remark}

\subsection{Supuestos de compartimentación}

\begin{description}
\item[A1. Mezcla homogénea.] Cada compartimento está uniformemente
distribuido en su volumen (Definición~\ref{def13}): un escalar describe su
estado y una medición puntual es representativa. Los gradientes locales en
la vecindad de la masa larval se aceptan como fuente de error de medición,
a vigilar en la validación.

\item[A2. Población homogénea.] Las $N(t)$ larvas comparten el estado
$(A, B, L)$ de una larva promedio (Definición~\ref{def31}); la mortalidad,
si ocurre, se admite mediante $N(t)$ decreciente.

\item[A3. Operación por lote.] Sin recargas de alimento ni de larvas en
$[0, t_f]$ (Observación~\ref{rem:batch}); el estado es continuo y las
únicas conmutaciones son las de ventilación (Definición~\ref{def42}).
\end{description}

\subsection{Supuestos fisiológicos}

\begin{description}
\item[B1. Asimilación limitante.] La asimilación de alimento es el proceso
metabólico limitante, $r_A = aB$ \citep{eriksenDynamicModellingFeed2022}, con $a$ modulada por
funciones logísticas del peso según el modelo de referencia.

\item[B2. Mantenimiento y costos de síntesis.] El mantenimiento es
proporcional a la biomasa estructural, $r_{\mathrm{CO_2},m} = mB$, y los
costos de síntesis $Y_B$, $Y_L$ son fracciones constantes de los flujos
respectivos (Ecs.~\eqref{eq:rco2m}--\eqref{eq:rco2l}).

\item[B3. Estado estacionario de asimilados.] El compartimento $A$
equilibra sus flujos en escalas rápidas frente a la dinámica de $B$ y $L$
(Ec.~1 de \citet{eriksenDynamicModellingFeed2022}), de modo que $\dot{A} \approx 0$ en
\eqref{eq:eriksen-edos}.

\item[B4. Parámetros constantes.] Se adoptan los parámetros fisiológicos
$(a, m, Y_B, Y_L, B_{\max}, \beta)$ como constantes durante el ciclo. Esta
es la práctica establecida en la literatura de referencia: el modelo de
\citet{eriksenDynamicModellingFeed2022} fue formulado y ajustado con parámetros constantes a
temperatura de cría controlada, con buen desempeño frente a datos de
distintos sustratos y humedades \citep{eriksenDynamicModellingFeed2022,eriksenMetabolicPerformanceFeed2024}. Dado que
el presente sistema mantiene $T$ regulada en la zona de confort larval
(Definición~\ref{def41}), la aproximación es razonable; la dependencia
explícita de los parámetros con $T$ (por ejemplo, factores de tipo
$Q_{10}$ o Arrhenius) queda planteada como extensión para el caso de
operación fuera de la zona regulada.
\end{description}
\section{Modelo matemático}\label{sec:modelo-general}

% 5.1 Din\'amica larval: forma general
% 5.2 Balance del sustrato
% 5.3 Din\'amica del lecho
% 5.4 Balances de gases en modo ventilado
% 5.5 Modo cerrado y forma unificada
% 5.6 Balance de humedad
% 5.7 Balance de energ\'ia y temperatura
% 5.8 Simplificaciones del modelo

Esta sección formula el modelo general del sistema, con todas las
dependencias explícitas. Las simplificaciones que se adoptarán para la
simulación y la estimación se enuncian al final de la sección
(Sección~\ref{sec:simplificaciones}), sobre el modelo ya construido.

\subsection{Din\'amica larval: forma general}
\label{sec:larval-general}

La din\'amica de la larva promedio sigue la estructura de compartimentos del modelo de Eriksen (Ecs.~\eqref{eq:eriksen-edos}), pero con las dependencias del sistema controlado hechas expl\'icitas. Como la larva vive y se alimenta dentro del lecho, las variables ambientales que modulan su fisiolog\'ia son las del lecho ---su temperatura $T_s$ y su contenido de agua $\theta_s$--- y no las del aire interno directamente. En su forma general, las tasas son funciones de
\begin{align}
r_A &= r_A(B, S, T_s, \theta_s; \theta) && \text{(asimilaci\'on)}, \label{eq:ra-general}\\
r_B &= r_B(A, B, T_s; \theta) && \text{(s\'intesis estructural)}, \label{eq:rb-general}\\
r_L &= r_L(A, B, L, T_s; \theta) && \text{(acumulaci\'on/reciclaje de l\'ipidos)}, \label{eq:rl-general}\\
r_{\mathrm{CO_2}} &= m(T_s)\,B + Y_B\,r_B + Y_L\,r_L && \text{(respiraci\'on total)}, \label{eq:rco2-general}
\end{align}
y el consumo de ox\'igeno se vincula a la respiraci\'on mediante el cociente respiratorio,
\begin{equation}
r_{\mathrm{O_2}} = \frac{r_{\mathrm{CO_2}}}{RQ}, \qquad RQ = RQ(A, B, L; \theta),
\label{eq:ro2-general}
\end{equation}
donde $RQ$ depende del sustrato metab\'olico que se oxida: pr\'oximo a $1.0$ cuando dominan carbohidratos de la dieta, y decreciente hacia $\approx 0.7$ cuando el reciclaje de l\'ipidos ($r_L < 0$) sostiene el metabolismo.

\begin{remark}[Acoplamiento a trav\'es del lecho]
\label{rem:dependencias}
Las dependencias escritas en~\eqref{eq:ra-general}--\eqref{eq:ro2-general} son deliberadamente generales. La asimilaci\'on depende del sustrato disponible $S$ (a sustrato agotado no hay asimilaci\'on), de la temperatura del lecho $T_s$ (la fisiolog\'ia es t\'ermicamente sensible) y de su contenido de agua $\theta_s$ (la actividad de las larvas en el lecho depende del agua disponible). N\'otese que $T_s$ y $\theta_s$ no son forzamientos ex\'ogenos como en el modelo original, sino estados del propio sistema (Definici\'on~\ref{def32}): el aire interno afecta a las larvas solo de manera indirecta, a trav\'es de la din\'amica t\'ermica e h\'idrica del lecho (Secci\'on~\ref{sec:lecho}), mientras que las larvas modifican el lecho con su calor, su respiraci\'on y su actividad, y el lecho, a su vez, intercambia con el aire. La din\'amica larval queda as\'i acoplada bidireccionalmente con el ambiente del contenedor, mediada por el lecho.
\end{remark}

\subsection{Balance del sustrato}
\label{sec:sustrato}

El sustrato es un compartimento activo (Observaci\'on~\ref{rem:microbiano}): lo consume la poblaci\'on larval y lo transforma la microbiota. Su balance general de masa es
\begin{equation}
\dot S = \underbrace{-\frac{N\,r_A}{\eta_A}}_{\text{consumo larval}} \;\underbrace{-\;k_{\mathrm{mic}}(T_s, \theta_s)\,S}_{\text{degradaci\'on microbiana}},
\label{eq:sustrato-general}
\end{equation}
donde $\eta_A \in (0,1]$ es la fracci\'on del sustrato ingerido que efectivamente se asimila (el resto pasa al residuo o frass, que permanece en el lote), $r_A$ es la tasa de asimilaci\'on por larva de la Ec.~\eqref{eq:ra-general} y $k_{\mathrm{mic}}(T_s, \theta_s)$ es una tasa espec\'ifica de degradaci\'on microbiana, dependiente de la temperatura y el contenido de agua del lecho.

Los procesos microbianos asociados son, en forma general,
\begin{align}
r^{\mathrm{mic}}_{\mathrm{CO_2}} &= r^{\mathrm{mic}}_{\mathrm{CO_2}}(S, T_s, \theta_s; \theta) && \text{(respiraci\'on microbiana)}, \label{eq:co2mic}\\
r^{\mathrm{mic}}_{\mathrm{O_2}} &= r^{\mathrm{mic}}_{\mathrm{O_2}}(S, T_s, \theta_s; \theta) && \text{(consumo microbiano de O$_2$)}, \label{eq:o2mic}\\
r_{\mathrm{CH_4}} &= r_{\mathrm{CH_4}}(S, T_s, \theta_s, \xi; \theta) && \text{(metanog\'enesis)}, \label{eq:ch4-general}
\end{align}
donde $\xi(t) \in [0,1]$ denota la fracci\'on anaerobia del lecho: la metanog\'enesis solo ocurre en micrositios sin ox\'igeno, y su intensidad depende de qu\'e tan desarrollada est\'e esa fracci\'on.

\begin{remark}[La fracci\'on anaerobia como variable latente]
\label{rem:xi}
La variable $\xi$ no es medible directamente ni es un estado con din\'amica propia en este nivel de descripci\'on: es una variable latente que crece con el contenido de agua del lecho $\theta_s$, la compactaci\'on y el consumo local de O$_2$, y decrece con la bioturbaci\'on larval. Esta lectura cualitativa tiene una consecuencia experimental ya anunciada (Observaci\'on~\ref{rem:tf}): al llegar la prepupa, cesa la bioturbaci\'on, $\xi$ aumenta y el CH$_4$ del outlet deber\'ia revelarlo. En la formulaci\'on simplificada (Secci\'on~\ref{sec:simplificaciones}) se decidir\'a si $\xi$ se parametriza como funci\'on de $(\theta_s, N, r_A)$ o se absorbe en los par\'ametros de~\eqref{eq:ch4-general}.
\end{remark}

\begin{remark}[El frass dentro del balance]
\label{rem:frass}
El residuo (frass) no aparece como compartimento: la fracci\'on no asimilada del sustrato permanece en el lote y sigue siendo parcialmente degradable por la microbiota, de modo que la distinci\'on operativa entre ``sustrato'' y ``frass'' es de calidad m\'as que de ubicaci\'on. En la forma general, \eqref{eq:sustrato-general} debe leerse como el balance de la materia org\'anica degradable del lecho; la partici\'on sustrato--frass es una simplificaci\'on pendiente (Secci\'on~\ref{sec:simplificaciones}).
\end{remark}

\subsection{Din\'amica del lecho}
\label{sec:lecho}

El lecho es la fase s\'olida del sistema: la mezcla de sustrato, larvas y frass que ocupa la base del contenedor. Su estado queda descrito por dos variables, el contenido de agua $\theta_s$ y la temperatura $T_s$, que gobiernan tanto la actividad microbiana del sustrato como, en gran parte, la actividad larval.

\paragraph{Balance de agua del lecho.} Sea $\theta_s$ el contenido gravim\'etrico de agua en base h\'umeda, de modo que la masa de agua del lecho es $M_{\mathrm{agua}} = \theta_s \cdot S$. En operaci\'on batch no hay entradas de agua al lecho, y adoptamos la forma m\'as simple del balance: la \'unica salida expl\'icita es la evaporaci\'on hacia el aire interno, $E_{\mathrm{lecho}}(\theta_s, T_s, w; \theta)$, en la que queda absorbido el consumo h\'idrico larval. La conservaci\'on de la masa de agua da
\begin{equation}
\frac{d}{dt}\left(\theta_s \cdot S\right) = -E_{\mathrm{lecho}}(\theta_s, T_s, w; \theta),
\label{eq:agua-lecho}
\end{equation}
y expandiendo la derivada, con $\dot S$ dado por~\eqref{eq:sustrato-general},
\begin{equation}
\dot\theta_s = -\frac{E_{\mathrm{lecho}}(\theta_s, T_s, w; \theta) + \theta_s \cdot \dot S}{S}.
\label{eq:theta-lecho}
\end{equation}
La condici\'on inicial $\theta_s(0) = \theta_{s,0}$ se determina gravim\'etricamente al momento de la carga (Secci\'on~6).

\paragraph{Nodo t\'ermico del lecho.} Toda la actividad metab\'olica del sistema ---larval y microbiana--- ocurre dentro del lecho, de modo que el calor metab\'olico se genera en este nodo y no en el aire. Con $C_{\mathrm{th},s}$ la capacidad t\'ermica del lecho, $h_s$ el coeficiente de intercambio lecho--aire a trav\'es del \'area efectiva $A_s$, y $\lambda$ el calor latente de vaporizaci\'on del agua,
\begin{equation}
C_{\mathrm{th},s}\,\dot T_s = \underbrace{\dot Q_{\mathrm{met}}}_{\text{calor metab\'olico}} \;-\; \underbrace{h_s A_s\,(T_s - T)}_{\text{intercambio con el aire}} \;-\; \underbrace{\lambda\,E_{\mathrm{lecho}}}_{\text{enfriamiento evaporativo}},
\label{eq:termica-lecho}
\end{equation}
donde la potencia metab\'olica es
\begin{equation}
\dot Q_{\mathrm{met}} = \gamma_q\left(N\,r_{\mathrm{CO_2}} + r^{\mathrm{mic}}_{\mathrm{CO_2}}\right),
\label{eq:calor-metabolico}
\end{equation}
con $\gamma_q$ el calor liberado por unidad de CO$_2$ producido (equivalente energ\'etico de la respiraci\'on, la regla de Thornton fija el calor por mol de O$_2$ consumido en $430$--$480$~kJ/mol~O$_2$~\citep{thorntonRelationOxygenHeat1917,hansenUseCalorespirometricRatios2004}; el calor por mol de CO$_2$ producido es ese valor dividido entre $RQ$, de modo que para oxidaci\'on de carbohidratos ($RQ=1$) $\gamma_q \approx 460$--$480$~kJ/mol~CO$_2$~\citep{lightonMeasuringMetabolicRates2008}, variable con el sustrato metab\'olico al igual que $RQ$ y $\gamma_w$. El t\'ermino latente $-\lambda E_{\mathrm{lecho}}$ contabiliza el enfriamiento evaporativo del lecho, que ocurre en este nodo; el calor de condensaci\'on, en cambio, se deposita en las paredes y se contabiliza en el nodo de aire (Secci\'on~\ref{sec:energia}).

\begin{remark}[Consistencia de la localizaci\'on de las fuentes]
\label{rem:lecho-consistencia}
El modelo de dos nodos corrige una inconsistencia de la formulaci\'on agregada: el calor, el CO$_2$ microbiano y la mayor parte del vapor de agua se generan en el lecho y llegan al aire por intercambio, no por generaci\'on directa en la fase gaseosa. Adem\'as, ambos estados del lecho son observables con el sensor capacitivo de suelo (con calibraci\'on gravim\'etrica para $\theta_s$), de modo que la extensi\'on del vector de estado de $10$ a $12$ variables no introduce estados inaccesibles.
\end{remark}

\subsection{Balances de gases en modo ventilado}
\label{sec:gases-ventilado}

Bajo el supuesto de mezcla homogénea del aire interno (A1), cada gas
$i \in \{\mathrm{CO_2}, \mathrm{O_2}, \mathrm{CH_4}\}$ obedece un balance de
moles en el volumen de aire $V_{\mathrm{aire}}$ del contenedor. En modo
ventilado ($Q > 0$):
\begin{align}
V_{\mathrm{aire}}\,\dot{c}_{\mathrm{CO_2}} &=
Q\,\bigl(c_{\mathrm{CO_2,in}} - c_{\mathrm{CO_2}}\bigr)
+ N\,r_{\mathrm{CO_2}} + r_{\mathrm{CO_2}}^{\mathrm{mic}},
\label{eq:co2-ventilado}\\
V_{\mathrm{aire}}\,\dot{c}_{\mathrm{O_2}} &=
Q\,\bigl(c_{\mathrm{O_2,in}} - c_{\mathrm{O_2}}\bigr)
- N\,r_{\mathrm{O_2}} - r_{\mathrm{O_2}}^{\mathrm{mic}},
\label{eq:o2-ventilado}\\
V_{\mathrm{aire}}\,\dot{c}_{\mathrm{CH_4}} &=
Q\,\bigl(c_{\mathrm{CH_4,in}} - c_{\mathrm{CH_4}}\bigr)
+ r_{\mathrm{CH_4}},
\label{eq:ch4-ventilado}
\end{align}
donde $c_i$ es la concentración del gas $i$ en el aire interno (igual a la
del \emph{outlet} por mezcla homogénea), $c_{i,\mathrm{in}}$ su
concentración en el aire de entrada, y las tasas $r$ corresponden a las
formas generales de las Ecs.~\eqref{eq:rco2-general},
\eqref{eq:ro2-general} y \eqref{eq:co2mic}--\eqref{eq:ch4-general}.
Nótese el signo del término larval y microbiano en cada ecuación:
producción en CO$_2$ y CH$_4$, consumo en O$_2$; y nótese que el CH$_4$
carece de término larval (B4 de la Sección~\ref{sec:supuestos}: metanogénesis solo en el
sustrato).

\begin{remark}[Estructura común de los balances]\label{rem:estructura-gas}
Las Ecs.~\eqref{eq:co2-ventilado}--\eqref{eq:ch4-ventilado} son instancias
del balance conservativo (Definición~\ref{def12}): el término en $Q$ es el
flujo que cruza la frontera y los términos en $r$ son la producción o
consumo interno. Dividiendo por $V_{\mathrm{aire}}$ se aprecia la constante
de tiempo del sistema de aire,
\begin{equation}
\tau_{\mathrm{aire}} = \frac{V_{\mathrm{aire}}}{Q},
\label{eq:tau-aire}
\end{equation}
que es el tiempo de residencia del aire en el contenedor. La separación de
escalas $\tau_{\mathrm{aire}} \ll t_f$ (minutos frente a días) implica que
en modo ventilado las concentraciones siguen casi instantáneamente a las
tasas de producción: en estado cuasi-estacionario,
\begin{equation}
c_i \approx c_{i,\mathrm{in}} \pm \frac{R_i}{Q},
\label{eq:cuasi-estacionario}
\end{equation}
con signo según producción o consumo. Esta aproximación es la que permite
estimar tasas a partir de mediciones de \emph{outlet} en operación normal,
y es complementaria del método de cámara estática~\citep{hutchinsonImprovedSoilCover1981,ermolaevGreenhouseGasEmissions2019}
(Ec.~\eqref{eq:modo-cerrado}).
\end{remark}

\begin{remark}[El caudal como entrada manipulada compartida]\label{rem:q-compartida}
$Q(t)$ aparece en las tres ecuaciones simultáneamente
(cf.~Observación~\ref{rem:jerarquia}): cualquier acción de control sobre
gases afecta también los balances de humedad y energía que siguen. En
\eqref{eq:cuasi-estacionario} se ve además el conflicto de diseño del lazo:
ventilar más diluye los gases (seguridad) pero reduce la señal $c_i -
c_{i,\mathrm{in}}$ disponible para estimar las tasas $R_i$ en operación
continua ---razón adicional para los cierres programados.
\end{remark}

\subsection{Modo cerrado y forma unificada}
\label{sec:modo-cerrado}

En modo cerrado ($Q = 0$ durante $[t_k^{c},\, t_k^{c}+\Delta]$,
Definición~\ref{def42}), los términos de flujo de las
Ecs.~\eqref{eq:co2-ventilado}--\eqref{eq:ch4-ventilado} se anulan y cada
concentración evoluciona solo por producción interna:
\begin{align}
V_{\mathrm{aire}}\,\dot{c}_{\mathrm{CO_2}} &=
N\,r_{\mathrm{CO_2}} + r_{\mathrm{CO_2}}^{\mathrm{mic}},
\label{eq:co2-cerrado}\\
V_{\mathrm{aire}}\,\dot{c}_{\mathrm{O_2}} &=
-\,N\,r_{\mathrm{O_2}} - r_{\mathrm{O_2}}^{\mathrm{mic}},
\label{eq:o2-cerrado}\\
V_{\mathrm{aire}}\,\dot{c}_{\mathrm{CH_4}} &= r_{\mathrm{CH_4}}.
\label{eq:ch4-cerrado}
\end{align}

Ambos modos se escriben de manera unificada usando la entrada conmutada
$Q(t)$ de la Ec.~\eqref{eq:q-conmutada}:
\begin{equation}
V_{\mathrm{aire}}\,\dot{c}_i =
Q(t)\,\bigl(c_{i,\mathrm{in}} - c_i\bigr) + \sigma_i\,R_i(t),
\qquad i \in \{\mathrm{CO_2}, \mathrm{O_2}, \mathrm{CH_4}\},
\label{eq:gas-unificado}
\end{equation}
donde $R_i$ es la tasa total (larval más microbiana cuando aplica) y
$\sigma_i = +1$ para gases producidos ($\mathrm{CO_2}$, $\mathrm{CH_4}$) y
$\sigma_i = -1$ para el consumido ($\mathrm{O_2}$).

\begin{remark}[Los dos modos como instrumentos de medición complementarios]
\label{rem:modos-complementarios}
La Ec.~\eqref{eq:gas-unificado} muestra que el sistema dispone de dos
regímenes de lectura de la misma cantidad $R_i(t)$. En modo ventilado,
cuasi-estacionario (Ec.~\eqref{eq:cuasi-estacionario}), $R_i$ se estima del
nivel: $R_i \approx Q\,(c_i - c_{i,\mathrm{in}})/\sigma_i$, con error
proporcional a la incertidumbre de $Q$ y de la diferencia de
concentraciones. En modo cerrado, $R_i$ se estima de la pendiente:
$R_i \approx V_{\mathrm{aire}}\,\dot{c}_i/\sigma_i$, independiente de $Q$ y
con señal creciente durante el cierre. La concordancia de ambas estimaciones
en los instantes previos y posteriores a cada cierre es una prueba de
consistencia interna del conjunto de sensores, análoga a la que aporta la
celda de carga (Ec.~\eqref{eq:celda-carga}) pero a nivel de gases
individuales.
\end{remark}

\begin{remark}[Transitorio de reapertura]\label{rem:reapertura}
Al finalizar un cierre, la reapertura de la ventilación restaura el modo
ventilado con condición inicial $c_i(t_k^{c}+\Delta)$ elevada (o deprimida,
en el caso del O$_2$), y la concentración relaja hacia el valor
cuasi-estacionario con la constante de tiempo $\tau_{\mathrm{aire}}$ de la
Ec.~\eqref{eq:tau-aire}. Este transitorio de relajación, típicamente de
pocos $\tau_{\mathrm{aire}}$, contiene información adicional sobre
$V_{\mathrm{aire}}$ efectivo y sobre la dinámica de mezcla, y no debe
descartarse como dato: forma parte del registro que el estimador puede
explotar (Sección~\ref{sec:mediciones}).
\end{remark}

\subsection{Balance de humedad}
\label{sec:humedad}

Sea $w(t)$ la humedad absoluta del aire interno (masa de vapor por unidad de masa de aire seco, o concentraci\'on molar de vapor, seg\'un la convenci\'on de unidades que se fije en la Secci\'on~2). Su balance en forma unificada es
\begin{equation}
V_{\mathrm{aire}}\,\dot w = Q(t)\,(w_{\mathrm{in}} - w) + E_{\mathrm{lecho}} + E_{\mathrm{met}} + J_{\mathrm{hum}}(t) - C_{\mathrm{cond}},
\label{eq:humedad-general}
\end{equation}
donde los cuatro t\'erminos internos son, en forma general:
\begin{align}
E_{\mathrm{lecho}} &= E_{\mathrm{lecho}}(\theta_s, T_s, w; \theta) && \text{(evaporaci\'on del lecho)}, \label{eq:evap-lecho}\\
E_{\mathrm{met}} &= \gamma_w\left(N\,r_{\mathrm{CO_2}} + r^{\mathrm{mic}}_{\mathrm{CO_2}}\right) && \text{(agua metab\'olica)}, \label{eq:agua-metabolica}\\
J_{\mathrm{hum}}(t) & && \text{(humidificaci\'on activa)}, \label{eq:humidificador}\\
C_{\mathrm{cond}} &= C_{\mathrm{cond}}(T, T_{\mathrm{pared}}, w) && \text{(condensaci\'on en paredes)}. \label{eq:condensacion}
\end{align}
En~\eqref{eq:agua-metabolica}, $\gamma_w$ es el agua producida por unidad de CO$_2$ respirado, fijada por la estequiometr\'ia de la oxidaci\'on del sustrato metab\'olico; su dependencia del tipo de sustrato oxidado es la misma que la del cociente respiratorio (Ec.~\eqref{eq:ro2-general}), pues carbohidratos, prote\'inas y l\'ipidos rinden distinta cantidad de agua por mol de CO$_2$. N\'otese que el agua metab\'olica es agua de \emph{nueva producci\'on}: no descuenta del stock h\'idrico del lecho en~\eqref{eq:agua-lecho}, sino que se libera al aire como producto de la respiraci\'on. El t\'ermino~\eqref{eq:humidificador} es la entrada manipulada del lazo de humedad (Definici\'on~\ref{def41}), y~\eqref{eq:condensacion} representa la remoci\'on de vapor por condensaci\'on cuando la temperatura de las paredes cae bajo el punto de roc\'io del aire interno.

\begin{remark}[La humedad como variable de acoplamiento indirecto]
\label{rem:humedad-doble}
$w$ es la variable del aire con m\'as interacciones del modelo: recibe la evaporaci\'on del lecho ---que ella misma modula a trav\'es de la demanda evaporativa $E_{\mathrm{lecho}}(\theta_s, T_s, w; \theta)$---, el agua metab\'olica y la humidificaci\'on activa, y su punto de roc\'io vincula~\eqref{eq:humedad-general} con el balance de energ\'ia a trav\'es de $T_{\mathrm{pared}}$. N\'otese, sin embargo, que la modulaci\'on de la actividad biol\'ogica por el agua ya no corre por cuenta de $w$ sino de $\theta_s$ (Ecs.~\eqref{eq:ra-general}, \eqref{eq:co2mic}--\eqref{eq:ch4-general}): el aire h\'umedo o seco afecta la fisiolog\'ia solo de manera indirecta, controlando la evaporaci\'on que seca o conserva el lecho. Estas interacciones son las que justifican el lazo de control de humedad como lazo independiente del de temperatura (Observaci\'on~\ref{rem:jerarquia}).
\end{remark}

\begin{remark}[Medici\'on: humedad relativa vs.\ absoluta]
\label{rem:hr-ha}
El sensor del sistema mide humedad relativa $H$, mientras que el balance~\eqref{eq:humedad-general} se escribe en humedad absoluta $w$, que es la variable extensiva (Definici\'on~\ref{def12}). La relaci\'on entre ambas pasa por la presi\'on de vapor de saturaci\'on, funci\'on fuertemente no lineal de $T$ (ecuaci\'on de Clausius--Clapeyron o aproximaciones tipo Antoine/Magnus): $H = p_v(w)\,/\,p_v^{\mathrm{sat}}(T)$. El mapa de observaci\'on $g$ de la Secci\'on~6 incluir\'a esta conversi\'on de forma expl\'icita ---un punto donde el modelo y la instrumentaci\'on deben hablar con cuidado, pues una lectura de $H$ constante no implica $w$ constante si $T$ var\'ia.
\end{remark}

\subsection{Balance de energ\'ia y temperatura}
\label{sec:energia}

La temperatura del aire interno $T$ no es una variable extensiva; su din\'amica se deriva del balance de energ\'ia sobre el aire del contenedor. Con el modelo de dos nodos de la Secci\'on~\ref{sec:lecho}, el calor metab\'olico se genera en el lecho (Ec.~\eqref{eq:calor-metabolico}) y llega al aire por intercambio; las p\'erdidas por paredes y la condensaci\'on se contabilizan contra el nodo de aire, que es la fase en contacto con ellas. Sea $C_{\mathrm{th}}$ la capacidad t\'ermica efectiva del aire interno (incluida la cara interna de las paredes, en J/K). En forma general,
\begin{equation}
C_{\mathrm{th}}\,\dot T = \underbrace{h_s A_s\,(T_s - T)}_{\text{intercambio con el lecho}} + \underbrace{P_{\mathrm{cal}}(t)}_{\text{calefactor}} - \underbrace{\dot Q_{\mathrm{vent}}}_{\text{ventilaci\'on}} - \underbrace{\dot Q_{\mathrm{pared}}}_{\text{p\'erdidas}} + \underbrace{\dot Q_{\mathrm{cond}}}_{\text{condensaci\'on}},
\label{eq:energia-general}
\end{equation}
con los t\'erminos dados por
\begin{align}
\dot Q_{\mathrm{vent}} &= Q(t)\,\rho_{\mathrm{aire}}\,c_{p,\mathrm{aire}}\,(T - T_{\mathrm{in}}) && \text{(arrastre sensible)}, \label{eq:calor-vent}\\
\dot Q_{\mathrm{pared}} &= U A_{\mathrm{pared}}\,(T - T_{\mathrm{amb}}) && \text{(conducci\'on por paredes)}, \label{eq:calor-pared}\\
\dot Q_{\mathrm{cond}} &= \lambda\,C_{\mathrm{cond}} && \text{(calor latente de condensaci\'on)}, \label{eq:calor-cond}
\end{align}
donde $\rho_{\mathrm{aire}} c_{p,\mathrm{aire}}$ es la capacidad calor\'ifica volum\'etrica del aire, $U A_{\mathrm{pared}}$ el coeficiente global de transferencia por las paredes hacia el ambiente exterior a $T_{\mathrm{amb}}$, y $\lambda$ el calor latente de vaporizaci\'on del agua. N\'otese la partici\'on de los t\'erminos latentes entre los dos nodos: el costo de evaporaci\'on $-\lambda E_{\mathrm{lecho}}$ lo paga el lecho (Ec.~\eqref{eq:termica-lecho}), donde ocurre la evaporaci\'on, y el calor de condensaci\'on $+\lambda C_{\mathrm{cond}}$ lo recibe el aire, donde se deposita el condensado; sumados sobre ambos nodos reproducen el t\'ermino neto $-\lambda\,(E_{\mathrm{lecho}} - C_{\mathrm{cond}})$ del balance agregado.

\begin{remark}[El calor metab\'olico como firma t\'ermica del crecimiento]
\label{rem:calor-firma}
La Ec.~\eqref{eq:calor-metabolico} liga directamente la din\'amica t\'ermica con la fisiolog\'ia: toda la entalp\'ia de los asimilados que no se almacena en biomasa se disipa como calor, y en la pr\'actica la potencia metab\'olica es proporcional a la tasa respiratoria~\citep{hansenUseCalorespirometricRatios2004}. En el modelo de dos nodos, esa firma se lee primero en el lecho: el exceso $T_s - T$ es una medici\'on indirecta de la actividad metab\'olica total ---complementaria del CO$_2$ y con distinta constante de tiempo (t\'ermica, no de mezcla)--- y el sensor de temperatura del lecho la registra sin intervenir el lote. El pico de $T$ sobre $T_{\mathrm{ref}}$ que el controlador debe remover es, a su vez, la sombra en el aire de ese pico de crecimiento.
\end{remark}

\begin{remark}[Las dos capacidades t\'ermicas]
\label{rem:cth}
El modelo de dos nodos reparte la capacidad t\'ermica del sistema en dos: la del lecho, $C_{\mathrm{th},s}$, y la del aire, $C_{\mathrm{th}}$. Ninguna es estrictamente constante durante el ciclo: la masa del lecho decrece (celda de carga) y su humedad var\'ia, de modo que en la forma general se admite $C_{\mathrm{th},s} = C_{\mathrm{th},s}(m_{\mathrm{LC}}, \theta_s)$; la del aire es mucho menor y m\'as estable. La decisi\'on de tomarlas constantes o dependientes del estado queda para la secci\'on de simplificaciones (Secci\'on~\ref{sec:simplificaciones}). Ambas son estimables experimentalmente de los transitorios t\'ermicos ante cambios de consigna del calefactor, sin intervenir el lote.
\end{remark}

\subsection{Simplificaciones del modelo}
\label{sec:simplificaciones}

\paragraph{El sistema acoplado en forma de estado.} Reuniendo los balances de las Secciones~\ref{sec:larval-general}--\ref{sec:energia}, el modelo general queda cerrado como un sistema din\'amico en el sentido de la Definici\'on~\ref{def11}.

\begin{definition}[Modelo general del sistema de bioconversi\'on]
\label{def:modelo-general}
El modelo general es el sistema $\dot x = f(x, u; \theta)$, $y = g(x; \theta)$, con estado
\begin{equation}
\mathbf{x} = (A, B, L, N, S, \theta_s, T_s, T, w, c_{\mathrm{CO_2}}, c_{\mathrm{O_2}}, c_{\mathrm{CH_4}}) \in \mathbb{R}^{12},
\label{eq:estado-12}
\end{equation}
entradas manipuladas $u = (Q, P_{\mathrm{cal}}, P_{\mathrm{hum}})$ m\'as las condiciones del inlet $(T_{\mathrm{in}}, w_{\mathrm{in}}, c_{i,\mathrm{in}})$, y din\'amica $f$ dada componente a componente por: los balances larvales~\eqref{eq:eriksen-edos} con las tasas generales~\eqref{eq:ra-general}--\eqref{eq:ro2-general}; el balance de sustrato~\eqref{eq:sustrato-general}; los balances del lecho~\eqref{eq:theta-lecho} y~\eqref{eq:termica-lecho}; el balance unificado de gases~\eqref{eq:gas-unificado}; el balance de humedad~\eqref{eq:humedad-general}; y el balance de energ\'ia del aire~\eqref{eq:energia-general}. El mapa de observaci\'on $g$ se desarrolla en la Secci\'on~6.
\end{definition}

Las formas funcionales de las tasas y los valores de los par\'ametros requieren decisiones de modelado que el sistema real debe justificar. Se adoptan las siguientes simplificaciones, todas ellas verificables o refinables con los datos del sistema instrumentado.

\paragraph{S1. Tasas larvales expl\'icitas con par\'ametros constantes.} Las tasas~\eqref{eq:ra-general}--\eqref{eq:rco2-general} toman las formas expl\'icitas del modelo de \citet{eriksenDynamicModellingFeed2022}: asimilaci\'on $r_A = a\,B$ con $a$ modulada por funciones log\'isticas del peso (estructura de estadios, $B_{\mathrm{m\'ax}}$, $\beta$), y costos $Y_B$, $Y_L$ constantes. La dependencia t\'ermica se congela: $m(T_s) = m$ constante, pues el control mantiene $T_s$ en la zona de confort larval (supuesto B4 de la Secci\'on~4). La dependencia con $\theta_s$ se admite como factor de disponibilidad del alimento solo si los datos lo exigen.

\paragraph{S2. Cin\'etica microbiana separable de primer orden.} La degradaci\'on microbiana se toma de primer orden en el sustrato, con tasa espec\'ifica separable,
\begin{equation}
k_{\mathrm{mic}}(T_s, \theta_s) = k_{\mathrm{mic}}^{\mathrm{ref}}\, f_T(T_s)\, f_\theta(\theta_s),
\label{eq:kmic-simplificada}
\end{equation}
con $f_T$ un factor t\'ermico tipo $Q_{10}$ o Arrhenius y $f_\theta$ un factor de humedad creciente y saturante (tipo Monod en el contenido de agua~\citep{monodGrowthBacterialCultures1949}). Las tasas de gases microbianos se ligan a la degradaci\'on por coeficientes de rendimiento constantes: $r^{\mathrm{mic}}_{\mathrm{CO_2}} = Y^{\mathrm{mic}}_{\mathrm{CO_2}}\, k_{\mathrm{mic}}\, S$ y $r^{\mathrm{mic}}_{\mathrm{O_2}} = r^{\mathrm{mic}}_{\mathrm{CO_2}} / RQ_{\mathrm{mic}}$. La alternativa Monod en $S$ (saturaci\'on por sustrato) queda como refinamiento: con una sola carga batch y $S$ abundante durante la mayor parte del ciclo, el primer orden es la hip\'otesis parsimoniosa.

\paragraph{S3. Parametrizaci\'on de la fracci\'on anaerobia.} La variable latente $\xi$ de~\eqref{eq:ch4-general} se parametriza como funci\'on est\'atica de los estados que la gobiernan,
\begin{equation}
\xi(\theta_s, N r_A) = \xi_{\mathrm{m\'ax}}\; \Lambda\!\left(\frac{\theta_s - \theta_s^{c}}{s_\theta}\right) \frac{1}{1 + \alpha_{\mathrm{biot}}\, N r_A},
\label{eq:xi-param}
\end{equation}
con $\sigma$ la funci\'on log\'istica, $\theta_s^{c}$ el umbral de humedad para anaerobiosis, y $\alpha_{\mathrm{biot}}$ el efecto de aireaci\'on por bioturbaci\'on larval, proporcional a la actividad de alimentaci\'on $N r_A$. La metanog\'enesis resulta $r_{\mathrm{CH_4}} = Y_{\mathrm{CH_4}}\, \xi\, k_{\mathrm{mic}}\, S$~\citep{batstoneIWAAnaerobicDigestion2002}. Esta forma reproduce cualitativamente la firma de terminaci\'on de la Observaci\'on~\ref{rem:tf}: al caer $r_A$ hacia la prepupa, $\xi$ sube y el CH$_4$ se dispara.

\paragraph{S4. Evaporaci\'on del lecho por demanda atmosf\'erica.} La evaporaci\'on se modela como flujo proporcional al d\'eficit de humedad del aire, limitado por la disponibilidad de agua en el lecho,
\begin{equation}
E_{\mathrm{lecho}} = k_e\, A_s \left( \varphi(\theta_s)\, w_{\mathrm{sat}}(T_s) - w \right)_+,
\label{eq:evap-simplificada}
\end{equation}
con $k_e$ un coeficiente de transferencia de masa, $w_{\mathrm{sat}}(T_s)$ la humedad de saturaci\'on a la temperatura del lecho (Clausius--Clapeyron/Magnus, cf.~Observaci\'on~\ref{rem:hr-ha}), $\varphi(\theta_s) \in [0,1]$ el factor de disponibilidad de agua, y $(\cdot)_+$ la parte positiva (no hay condensaci\'on sobre el lecho en este nivel de descripci\'on).

\paragraph{S5. Capacidades t\'ermicas constantes.} Se toman $C_{\mathrm{th}}$ y $C_{\mathrm{th},s}$ constantes durante el ciclo (cf.~Observaci\'on~\ref{rem:cth}), con $C_{\mathrm{th},s}$ estimada de la composici\'on inicial del lecho. La dependencia $C_{\mathrm{th},s}(m_{\mathrm{LC}}, \theta_s)$ queda como refinamiento si los transitorios t\'ermicos de calibraci\'on muestran deriva sistem\'atica.

\paragraph{S6. Coeficientes estequiom\'etricos constantes.} Se fijan $RQ$, $\gamma_q$ y $\gamma_w$ en valores constantes representativos de la dieta (mezcla carbohidratos--prote\'inas: $RQ \approx 0.9$--$1.0$), y $\gamma_q \approx 470$~kJ/mol CO$_2$~\citep{hansenUseCalorespirometricRatios2004,lightonMeasuringMetabolicRates2008}. La variaci\'on de $RQ$ con el sustrato metab\'olico (Ec.~\eqref{eq:ro2-general}) se usar\'a solo como diagn\'ostico sobre los datos, no dentro de la din\'amica simulada.

\begin{remark}[El caso degenerado de un solo nodo t\'ermico]
\label{rem:un-nodo}
Si el intercambio lecho--aire es r\'apido frente a las dem\'as din\'amicas ($h_s A_s \to \infty$), las Ecs.~\eqref{eq:termica-lecho} y~\eqref{eq:energia-general} colapsan a $T_s = T$ y el sistema se reduce a diez estados, con un \'unico balance t\'ermico agregado. No se adopta esa reducci\'on aqu\'i por dos razones: el sensor de suelo mide $T_s$ y $\theta_s$ directamente, de modo que ambos estados son observables sin costo adicional, y la diferencia $T_s - T$ es precisamente la firma t\'ermica de la actividad metab\'olica (Observaci\'on~\ref{rem:calor-firma}). El caso de un nodo queda, sin embargo, disponible como hip\'otesis nula en la validaci\'on: si los datos muestran $T_s \approx T$ en todo el ciclo, el modelo se simplifica solo.
\end{remark}

\begin{remark}[Econom\'ia del modelo resultante]
\label{rem:economia-modelo}
Las simplificaciones S1--S6 reducen el modelo general a un sistema r\'igido pero tractable: doce estados, tres entradas manipuladas y un vector de par\'ametros $\theta$ que agrupa los fisiol\'ogicos de Eriksen $(a, m, Y_B, Y_L, B_{\mathrm{m\'ax}}, \beta, \eta_A)$, los microbianos $(k_{\mathrm{mic}}^{\mathrm{ref}}, Y^{\mathrm{mic}}_{\mathrm{CO_2}}, RQ_{\mathrm{mic}}, Y_{\mathrm{CH_4}}, \xi_{\mathrm{m\'ax}}, \theta_s^c, s_\theta, \alpha_{\mathrm{biot}})$, y los f\'isicos $(k_e, h_s A_s, U A_{\mathrm{pared}}, C_{\mathrm{th}}, C_{\mathrm{th},s}, V_{\mathrm{aire}})$. La estimabilidad de cada grupo se discute en la Secci\'on~6: los cierres de ventilaci\'on identifican las tasas de gases, los transitorios de consigna los par\'ametros t\'ermicos, y el lote testigo sin larvas separa el bloque microbiano del larval.
\end{remark}
\section{Diseño del ambiente controlado}\label{sec:diseno}

Esta sección especifica el montaje físico del sistema experimental: la cámara,
la línea neumática de ventilación, los sensores, los actuadores y la cadena de
adquisición. El objetivo es que cada término del modelo de la
Sección~\ref{sec:modelo-general} tenga un elemento de hardware asociado: cada
entrada $u$ un actuador, cada componente de $y = g(x)$ un sensor, y cada
supuesto estructural (Sección~\ref{sec:supuestos}) una decisión de diseño
verificable.

\subsection{Elementos generales del diseño}\label{sec:elementos-generales}

\subsubsection{Arquitectura}

El montaje se organiza en seis subsistemas:

\begin{description}
  \item[Cámara.] Hielera rígida de polietileno de $5$~L de capacidad nominal
  (Truper HIEL-5, $26 \times 19 \times 18$~cm; volumen interior
  $\approx 4{,}5$~L), que combina aislamiento térmico, cierre hermético y bajo
  costo. En su interior se aloja el lecho de bioconversión.
  \item[Línea neumática.] Circuito de ventilación con entrada (\emph{inlet}) y
  salida (\emph{outlet}) independientes en la tapa, que impone el flujo $Q(t)$
  del balance de gases y permite conmutar entre ventilación y cierre estático
  (método de cámara cerrada).
  \item[Sensores.] Instrumentación de gases, temperatura, humedad y masa
  (Sección~\ref{sec:sensores}).
  \item[Actuadores.] Celda Peltier con disipador, humidificador ultrasónico y
  ventilación (Sección~\ref{sec:actuadores}).
    \item[Adquisición y control.] Microcontroladores ESP32 dedicados al
  monitoreo y al control: muestrean los sensores, ejecutan los lazos sobre
  los actuadores y transmiten las mediciones en línea a una base de datos
  en la nube (Sección~\ref{sec:dashboard}).
  \item[Dashboard.] Tablero web definido en la Sección~\ref{sec:dashboard},
  donde se especifican sus componentes mínimos; despliega las series en
  tiempo real y alimenta el gemelo digital del modelo de la
  Sección~\ref{sec:modelo-general}.  
  \item[Gemelo digital.] Instancia en línea del modelo de la
  Sección~\ref{sec:modelo-general}, alimentada por las mediciones.
\end{description}

\begin{remark}[Calefacción y enfriamiento unificados]\label{rem:peltier}
Una celda Peltier es un dispositivo termoeléctrico reversible: el signo de la
corriente determina si bombea calor hacia adentro o hacia afuera de la cámara.
El vector de entradas del modelo queda entonces
\[
  u = \big(Q,\; P_{\mathrm{pel}},\; P_{\mathrm{hum}}\big),
  \qquad
  P_{\mathrm{pel}} \in [-P_{\max},\, +P_{\max}],
\]
donde $P_{\mathrm{pel}}$ con signo reemplaza a la potencia de calefacción
$P_{\mathrm{cal}}$: un solo actuador cubre ambos modos, lo que simplifica el
hardware y el modelo.
\end{remark}

% fig:unidad-controlada se mantiene aquí (no la borres)

\subsubsection{Dimensionamiento del lote y del volumen de aire}

El experimento consiste en un único ciclo de bioconversión de duración
$t_f \in [14, 20]$~días, con una carga inicial de $N_0 \in [200, 700]$ larvas
y alimento en proporción $1{,}4$~g por larva. La densidad superficial de
cría es de $\approx 2{,}1$ larvas/cm$^2$.

\begin{table}[htbp]
  \centering
  \caption{Escenarios de carga del lote (área del lecho $A_s \approx 330$~cm$^2$).}
  \label{tab:carga-lote}
  \begin{tabular}{cccc}
    \hline
    $N_0$ (larvas) & $S_0$ (g) & Espesor del lecho & $V_{\mathrm{aire}}$ \\
    \hline
    200 & 280  & $\approx 1{,}7$~cm & $\approx 3{,}0$~L \\
    400 & 560  & $\approx 3{,}5$~cm & $\approx 2{,}8$~L \\
    700 & 980  & $\approx 6{,}0$~cm & $\approx 2{,}5$~L \\
    \hline
  \end{tabular}
\end{table}

El volumen de aire determina la duración de los cierres estáticos: con las
tasas de CO$_2$ por larva de la Sección~\ref{sec:caracterizacion}, la ventana
$\Delta$ admisible por fase es la de la Tabla~\ref{tab:delta-adaptativo}.

\begin{table}[htbp]
  \centering
  \caption{Ventana de cierre $\Delta$ por fase fisiológica
  ($N = 400$, $V_{\mathrm{aire}} \approx 3$~L; a $N = 700$ los tiempos se
  reducen en el factor $4/7$).}
  \label{tab:delta-adaptativo}
  \begin{tabular}{lcc}
    \hline
    Fase & Pendiente esperada (ppm/min) & $\Delta$ sugerida \\
    \hline
    Neonatos      & $\sim 10^2$  & $60$--$120$~min \\
    Crecimiento   & $10^2$--$10^3$ & $6$--$30$~min \\
    Pico metabólico & $\gtrsim 2 \times 10^3$ & $1{,}5$--$5$~min \\
    Prepupa       & $\sim 10^3$  & $9$--$18$~min \\
    \hline
  \end{tabular}
\end{table}

\begin{remark}[$\Delta$ operativo]\label{rem:delta-operativo}
$\Delta$ se acota por abajo por la resolución de los sensores
(ecuación~\eqref{eq:delta-min}) y por arriba por el piso de O$_2$ tolerable por
las larvas (ecuación~\eqref{eq:delta-max}). En el pico metabólico con
$N = 700$, $\Delta$ baja a $\approx 1{,}5$~min, cierres muy cortos que exigen
muestreo rápido; se acepta este régimen porque dura pocos días.
\end{remark}

\subsubsection{Línea neumática: \emph{inlet} y \emph{outlet}}

La línea conecta con la cámara mediante dos racores pasamuros herméticos
(conectores \emph{push-to-connect} de $1/8''$): el \emph{inlet} se ubica en
la parte inferior de una cara lateral y el \emph{outlet} en la tapa o en la
parte superior de una cara lateral. La posición definitiva de ambos puertos
se definirá mediante una simulación CFD del flujo en el volumen de aire
(Ansys o FEniCS/\texttt{dolfinx}), con el fin de garantizar el supuesto de
mezclado homogéneo de la Sección~\ref{sec:supuestos} y evitar zonas muertas
sobre el lecho. El diámetro de toda la línea queda fijado por estos racores:

\begin{description}
  \item[Tubería.] Poliuretano $6 \times 4$~mm (exterior $\times$ interior),
  tramos de longitud total menor a $1$~m. Con el flujo máximo de diseño
  $Q_{\max} \approx 1{,}7$~L/min la velocidad en el tubo es
  $v = Q/A \approx 2{,}3$~m/s y la caída de presión es despreciable.
  \item[Bomba de entrada.] Microbomba de diafragma de $12$~V DC, caudal
  nominal $2$~L/min, accionada por PWM desde el ESP32; es el actuador que
  impone $Q(t)$.
  \item[Bomba de salida.] Segunda bomba idéntica, operada con el mismo duty
  cycle. Usar bombas gemelas en entrada y salida mantiene la presión interior
  cercana a la ambiente: la cámara es rígida y una sola bomba la
  sobrepresurizaría o despresurizaría, sesgando la lectura del BMP390 y la
  estimación de concentraciones. Como respaldo, el \emph{outlet} lleva un
  filtro hidrofóbico de $0{,}2$~$\mu$m que evita escapes de aerosoles durante
  la ventilación.
  \item[Estrategia de cierre.] Ambas bombas se detienen simultáneamente y el
  circuito queda sellado por las propias bombas de diafragma (válvulas de
  retención internas); no se requieren electroválvulas adicionales.
\end{description}

El caudal se dimensiona para renovar el volumen de aire en el régimen más
exigente (pico metabólico, $\Delta \approx 1{,}5$~min):
\begin{equation}\label{eq:q-dimensionamiento}
  Q_{\mathrm{dis}}
  \;\geq\; \frac{V_{\mathrm{aire}}}{3\,\Delta_{\min}}
  \;\approx\; \frac{2{,}5~\mathrm{L}}{3 \times 1{,}5~\mathrm{min}}
  \;\approx\; 0{,}6~\mathrm{L/min},
\end{equation}
y se elige $Q_{\mathrm{dis}} = 1$~L/min nominal (bomba de $2$~L/min al
$50\%$), con margen hasta $1{,}7$~L/min para el escenario $N = 700$. El
tiempo característico de renovación es
$\tau_{\mathrm{aire}} = V_{\mathrm{aire}}/Q \approx 2{,}5$~min, coherente con
la ventilación conmutada de la ecuación~\eqref{eq:q-conmutada}.

\begin{remark}[Muestreo durante los cierres]\label{rem:muestreo-cierres}
Con bombas detenidas el NDIR y el SHT35 muestrean cada $10$~s; un cierre de
$\Delta = 1{,}5$~min aporta $\sim 9$ puntos por tramo, suficiente para ajustar
la pendiente por mínimos cuadrados.
\end{remark}

\begin{remark}[Escalado y réplicas]\label{rem:escalado}
El diseño se valida con un único experimento de $14$ a $20$~días; las réplicas
se ejecutan de forma secuencial al finalizar el primer ciclo, reutilizando el
mismo montaje tras la limpieza y recalibración en vacío
(Sección~\ref{sec:caracterizacion}). No se requiere hardware adicional: la
replicación es un problema de calendario, no de equipamiento.
\end{remark}
\section{Mediciones y estimación}\label{sec:mediciones}

El modelo de la Sección~\ref{sec:modelo-general} solo es útil si sus
estados y parámetros pueden recuperarse de las mediciones disponibles.
Esta sección define el mapa de observación del sistema instrumentado
(Sección~\ref{sec:diseno}), analiza qué es observable e identificable
con esa instrumentación, y establece el protocolo de caracterización
previa y la estrategia de estimación que alimenta el gemelo digital.

\subsection{Vector de mediciones y mapa de observación}\label{sec:mapa-obs}

La instrumentación del diseño aporta tres grupos de señales, muestreadas
a $0{,}1$~Hz ($10$~s entre muestras, Observación~\ref{rem:muestreo-cierres})
y registradas con marca temporal común:

\begin{description}
  \item[Gases y aire interno (outlet).] Temperatura $T$ y humedad relativa
  $H$ (SHT35), CO$_2$ (NDIR), CH$_4$ y O$_2$ en la corriente de salida:
  \[
    \mathbf{y}_{\mathrm{gas}}(t) =
    \bigl(T,\; H,\; c_{\mathrm{CO_2}},\; c_{\mathrm{CH_4}},\;
    c_{\mathrm{O_2}}\bigr)(t).
  \]
  Bajo el supuesto de mezcla homogénea (A1), estas lecturas son los
  propios estados del aire interno: el mapa de observación es aquí la
  identidad, más ruido de medición.
  \item[Lecho.] El sensor capacitivo insertado en el lecho aporta
  $(\theta_s^{\mathrm{sen}}, T_s^{\mathrm{sen}})$. La temperatura lee el
  estado $T_s$ directamente; el canal capacitivo requiere calibración
  gravimétrica para relacionarse con $\theta_s$
  (Sección~\ref{sec:caracterizacion}):
  \[
    \theta_s^{\mathrm{sen}}(t) = h_\theta\bigl(\theta_s(t)\bigr) + \varepsilon_\theta,
    \qquad
    T_s^{\mathrm{sen}}(t) = T_s(t) + \varepsilon_T.
  \]
  \item[Masa del lote.] La celda de carga registra $m_{\mathrm{LC}}(t)$,
  la masa total del lote, sujeta a la restricción de consistencia global
  de la Ec.~\eqref{eq:celda-carga}: su tasa de decrecimiento debe
  igualar el flujo neto de masa por el \emph{outlet}.
\end{description}

Además, las variables de supervisión (H$_2$S, NH$_3$) se registran con
fines de alarma sin entrar al mapa de observación del modelo
(Sección~\ref{sec:sistema}), y las condiciones del \emph{inlet}
$(T_{\mathrm{in}}, w_{\mathrm{in}}, c_{\mathrm{O_2,in}})$ se miden como
forzamientos conocidos, no como estados.

\subsection{Observabilidad con $(T, H, \mathrm{CO_2}, \mathrm{CH_4}, \mathrm{O_2})$}\label{sec:observabilidad}

El vector de estado tiene dimensión $12$
(Ec.~\eqref{eq:estado-completo}). Su observabilidad se analiza por
grupos de compartimentos:

\begin{description}
  \item[Aire interno $(T, w, c_{\mathrm{CO_2}}, c_{\mathrm{CH_4}}, c_{\mathrm{O_2}})$.]
  Observable directamente: cinco estados medidos por identidad, con $w$
  obtenido de $H$ y $T$ vía psicrometría.
  \item[Lecho $(\theta_s, T_s)$.] Observable con el sensor capacitivo
  calibrado y el balance de agua~\eqref{eq:agua-lecho} como verificación
  cruzada.
  \item[Sustrato $S$.] No medible en línea de forma directa; se estima
  por el balance~\eqref{eq:sustrato-general} una vez conocidas las tasas,
  con ancla en la celda de carga y en la determinación gravimétrica
  inicial.
  \item[Larva promedio $(A, B, L)$ y población $N$.] Estados latentes:
  ningún sensor los lee. Su única vía de observación son las tasas
  agregadas $R_i$ medidas en los cierres
  (Observación~\ref{rem:cierres}), que por~\eqref{eq:rco2-general}
  combinan $mB$, $Y_B r_B$ y $Y_L r_L$. El par $(\mathrm{CO_2}, \mathrm{O_2})$
  aporta además el cociente respiratorio medido~\eqref{eq:rq-medido},
  que discrimina el sustrato metabólico oxidado.
\end{description}

En síntesis: siete estados son medibles de forma directa o semidirecta,
y cinco $(A, B, L, N, S)$ son estimables solo a través del modelo. Esta
estructura justifica la estrategia de estimación de la
Sección~\ref{sec:estimacion}: los estados medidos anclan, y las tasas de
los cierres informan a los latentes.

\subsection{Identificabilidad: separación larva--microbiota}\label{sec:identificabilidad}

Las señales del \emph{outlet} son superposiciones
(Ec.~\eqref{eq:co2-superpuesto}): sin información adicional, la
contribución larval y la microbiana no son separables, y el
$RQ_{\mathrm{medido}}$ no es interpretable fisiológicamente
(Ec.~\eqref{eq:rq-medido}). El protocolo resuelve esta identificabilidad
con tres mecanismos complementarios:

\begin{enumerate}
  \item \textbf{Línea base microbiana temprana.} El primer cierre de
  ventilación, ejecutado dentro de las primeras $24$~h tras la carga,
  mide una señal dominada por $r^{\mathrm{mic}}$ (larvas neonatas de
  masa despreciable). Es la estimación inicial de la actividad
  microbiana del sustrato.
  \item \textbf{Lote testigo sin larvas.} Una corrida completa con el
  mismo sustrato, humedad inicial y calendario de cierres, pero
  $N_0 = 0$, mide $r^{\mathrm{mic}}_{\mathrm{CO_2}}$ y
  $r^{\mathrm{mic}}_{\mathrm{O_2}}$ de forma independiente en todo el
  horizonte. Este tratamiento es parte obligatoria del protocolo
  experimental: sin él, la descomposición
  de~\eqref{eq:co2-superpuesto} queda indeterminada.
  \item \textbf{Cierre post-cosecha.} Un cierre adicional tras retirar
  las prepupas mide la respiración residual del frass, que acota la
  actividad microbiana al final del ciclo.
\end{enumerate}

\begin{remark}[El CH$_4$ como testigo de anaerobiosis]
La metanogénesis es casi exclusivamente microbiana
(Observación~\ref{rem:microbiano}), de modo que el CH$_4$ no sufre el
problema de superposición: su señal informa directamente la fracción
anaerobia latente $\xi$ (Observación~\ref{rem:xi}) y valida la firma de
terminación propuesta en la Observación~\ref{rem:tf}.
\end{remark}

\subsection{Caracterización previa del sistema}\label{sec:caracterizacion}

Antes de cargar el primer lote, el sistema vacío y los sensores se
caracterizan en cuatro pruebas, cuyos resultados son parámetros del
modelo o de la cadena de medición:

\begin{enumerate}
  \item \textbf{Prueba de hermeticidad} (Observación~\ref{rem:fugas}):
  estimación de $k_{\mathrm{fuga}}$ por decaimiento de CO$_2$ inyectado.
  Si $k_{\mathrm{fuga}}$ no es despreciable, se incorpora como término
  correctivo en los balances de gases.
  \item \textbf{Volumen de aire efectivo.} $V_{\mathrm{aire}}$ se mide
  por dilución: inyección de un pulso conocido de CO$_2$ en la cámara
  cerrada con el lecho simulado, y lectura del incremento de
  concentración en estado estacionario. Este valor, no el geométrico,
  es el que entra en~\eqref{eq:modo-cerrado}.
  \item \textbf{Calibración de sensores.} El NDIR de CO$_2$ se calibra
  en dos puntos (aire ambiente $\sim 420$~ppm y gas patrón); el sensor
  capacitivo del lecho se calibra contra determinaciones gravimétricas
  de $\theta_s$ en el rango de operación.
  \item \textbf{Validación contra cromatografía.} En cada corrida se
  toman muestras por el septum en instantes seleccionados (inicio de
  cierre, mitad y fin) para análisis por cromatografía de gases. La
  comparación GC--sensor en línea cuantifica el sesgo y la deriva de
  los sensores de CO$_2$ y CH$_4$, y documenta la calidad metrológica
  del experimento.
\end{enumerate}

Las mismas pruebas se repiten en la recalibración entre réplicas
secuenciales (Observación~\ref{rem:escalado}).

\subsection{Estimación de estados y parámetros}\label{sec:estimacion}

La estimación se organiza en dos niveles, coherentes con la distinción
entre datos de identificación y datos de validación
(Observación~\ref{rem:id-lazo}):

\paragraph{Nivel 1: tasas instantáneas.} En cada cierre
$[t_k^c, t_k^c + \Delta]$, las pendientes de acumulación se ajustan por
mínimos cuadrados sobre los $\sim 9$ puntos disponibles
(Observación~\ref{rem:muestreo-cierres}) y entregan las tasas agregadas
\[
  R_i(t_k) = \sigma_i\, V_{\mathrm{aire}}\,
  \widehat{\left.\dot c_i\right|_{[t_k^c,\, t_k^c + \Delta]}},
  \qquad i \in \{\mathrm{CO_2}, \mathrm{O_2}, \mathrm{CH_4}\},
\]
independientes de $Q$ y sin modelado fisiológico. La ventana $\Delta$ se
ajusta día a día según las cotas~\eqref{eq:delta-min}
y~\eqref{eq:delta-max}, evaluadas con las tasas del día anterior: es la
primera función del gemelo digital.

\paragraph{Nivel 2: estados y parámetros del modelo.} Las series
$R_i(t_k)$, las señales continuas de la Sección~\ref{sec:mapa-obs} y las
condiciones iniciales conocidas alimentan el ajuste de los parámetros
$\boldsymbol{\theta}$ (fisiológicos, microbianos, térmicos e hídricos)
y la reconstrucción de los estados latentes $(A, B, L, N, S)$. La
vía prevista es una red neuronal informada por física (PINN), con las
ecuaciones de la Sección~\ref{sec:modelo-general} como restricciones de
la pérdida; un filtro de Kalman extendido o de ensamble sobre la forma
unificada ventilado--cerrado sirve como referencia comparativa. Los
datos de operación regulada normal se reservan para validación: el
modelo ajustado debe predecir las concentraciones entre cierres sin
haberlas usado en el ajuste.

\paragraph{Criterio de éxito.} La estimación se considera exitosa si
(i) los residuos de validación de $c_{\mathrm{CO_2}}$ y
$c_{\mathrm{O_2}}$ son compatibles con el ruido de los sensores,
(ii) la restricción de la celda de carga~\eqref{eq:celda-carga} se
cumple dentro de su tolerancia, y (iii) los parámetros fisiológicos
estimados caen en los rangos reportados para el modelo de referencia
\citep{eriksenDynamicModellingFeed2022,eriksenMetabolicPerformanceFeed2024}.

\subsection{Dashboard del sistema monitoreado}\label{sec:dashboard}

El tablero web es la interfaz del gemelo digital y despliega, como
mínimo, cuatro componentes:

\begin{enumerate}
  \item \textbf{Series en tiempo real} de $\mathbf{y}_{\mathrm{gas}}$,
  $(\theta_s^{\mathrm{sen}}, T_s^{\mathrm{sen}})$ y $m_{\mathrm{LC}}$,
  con los eventos de cierre marcados sobre las series de gases.
  \item \textbf{Tasas estimadas} $R_i(t_k)$ del Nivel 1, superpuestas
  con la predicción del modelo ajustado (Nivel 2): la brecha entre
  ambas es el diagnóstico en línea del gemelo.
  \item \textbf{Indicadores fisiológicos derivados:}
  $RQ_{\mathrm{medido}}$ y la razón CH$_4$/CO$_2$, con la firma de
  terminación de la Observación~\ref{rem:tf} evaluada en tiempo real.
  \item \textbf{Supervisión y alarmas}: H$_2$S y NH$_3$ contra umbrales,
  estado de los actuadores $(Q, P_{\mathrm{pel}}, P_{\mathrm{hum}})$ y
  cumplimiento de las cotas de seguridad de O$_2$ durante los cierres.
\end{enumerate}

La persistencia es una base de datos en la nube alimentada por los
ESP32 (Sección~\ref{sec:diseno}); cada réplica queda archivada como un
conjunto completo de series, eventos y metadatos del lote (cepa, edad y
cantidad de neonatos, $\theta_{s,0}$, $S_0$), reutilizable en el
análisis estadístico posterior.

\section{Discusión}\label{sec:discusion}

\subsection{Contribuciones del trabajo}

El documento aporta cuatro resultados que, en conjunto, no existen en la
literatura de bioconversión con \textit{H.~illucens}:

\begin{enumerate}
  \item \textbf{Un modelo unificado alimento--larvas--ambiente.} El modelo
  de \citet{eriksenDynamicModellingFeed2022} describe la larva en un ambiente
  exógeno; aquí la temperatura, la humedad y los gases del contenedor pasan
  a ser estados acoplados (Ec.~\eqref{eq:estado-completo}), con el lecho como nodo hidrotérmico
  intermedio y la forma unificada ventilado--cerrado de la
  Ec.~\eqref{eq:gas-unificado} como instrumento de medición.
  \item \textbf{Un método de medición híbrido con ventana adaptativa.}
  Los cierres de ventilación se dimensionan desde el propio modelo
  (Ecs.~\eqref{eq:delta-min}--\eqref{eq:delta-max},
  Tabla~\ref{tab:delta-adaptativo}): la literatura usa cierres fijos
  \citep{ermolaevGreenhouseGasEmissions2019} o muestreo continuo sin
  acumulación \citep{parodiBioconversionEfficienciesGreenhouse2020a}.
  \item \textbf{Un diseño experimental que resuelve la identificabilidad
  larva--microbiota.} El testigo sin larvas, la línea base temprana y el
  cierre post-cosecha (Sección~\ref{sec:identificabilidad}) separan las dos
  fuentes de CO$_2$ por diseño, no por suposición; la celda de carga añade
  una restricción de consistencia global independiente
  (Ec.~\eqref{eq:celda-carga}).
  \item \textbf{Un gemelo digital con estimación en línea.} La firma de
  terminación CH$_4$/CO$_2$ (Observación~\ref{rem:tf}), el ajuste adaptativo
  de $\Delta$ y la reconstrucción de estados latentes por PINN
  (Sección~\ref{sec:pinn}) convierten el montaje en un sistema de estimación,
  no solo de registro.
\end{enumerate}

\subsection{Posición frente a la literatura}

\begin{table}[htbp]
  \centering
  \caption{Posición del trabajo frente a los estudios de referencia.}
  \label{tab:posicion}
  \begin{tabular}{llll}
    \hline
    Estudio & Medición & Modelo & Estimación en línea \\
    \hline
    Eriksen (2022) & respirometría de larvas & larva, ambiente exógeno & no \\
    Ermolaev (2019) & cámara cerrada + GC puntual & no & no \\
    Parodi (2020) & cámara abierta continua & balance de masa estático & no \\
    Rossi (2024) & 4 cámaras abiertas, 12 h & interpolación empírica & no \\
    Este trabajo & conmutada, $\Delta$ adaptativo, GC & compartimental acoplado & sí (PINN) \\
    \hline
  \end{tabular}
\end{table}

La Tabla~\ref{tab:posicion} resume la posición: cada estudio aporta una
pieza ---el modelo fisiológico, la validación cromatográfica, el balance de
masa, la dinámica temporal--- y el presente trabajo los integra en un solo
marco donde el modelo y la medición se diseñan uno para el otro.

\subsection{Limitaciones}

\begin{description}
  \item[Réplicas secuenciales.] La confusión corrida--tiempo es inherente al
  equipamiento único (Observación~\ref{rem:paralelizacion}); la variabilidad
  entre lotes queda confundida con la variabilidad temporal del laboratorio,
  y solo se mitiga con metadatos, recalibración y el testigo común.
  \item[Ceguera parcial entre cierres.] La delegación del CH$_4$ al GC y del
  O$_2$ fino al modelo (Observación~\ref{rem:o2-supervision}) implica que
  entre cierres esas dos variables no se observan: un evento anaerobio
  breve entre cierres solo lo vería la alarma MOS.
  \item[Mezcla homogénea y gradientes del lecho.] El supuesto A1 se degrada
  en el escenario $N_0 = 700$ (lecho de $\approx 6$~cm), donde la lectura
  puntual del fondo puede subestimar la temperatura de la zona activa
  (Observación sobre lectura puntual, Sección~\ref{sec:camara-lecho}).
  \item[Población homogénea.] El supuesto A2 ignora el ensanchamiento de la
  distribución de pesos (Observación~\ref{rem:homogeneidad}); las réplicas
  acotan su efecto, pero no lo eliminan.
  \item[Una sola condición.] E5 prueba extrapolación en densidad, no en
  sustrato ni en temperatura: la generalización del modelo a otros residuos
  queda por verificar.
\end{description}

\subsection{Una lección de verificación de fuentes}

Durante la elaboración de este documento se detectó que un estudio antes
citado como respaldo metodológico en este campo ---sobre el efecto de la
humedad en las emisiones durante la bioconversión--- fue retractado por su
editorial. La referencia fue sustituida por un estudio equivalente del grupo
de SLU \citep{lindbergProcessEfficiencyGreenhouse2022}. El episodio deja
una regla metodológica que se adopta en lo sucesivo: toda referencia que
sustente una decisión experimental se verifica antes de citarse (estado de
retracción, existencia de los datos reportados), práctica especialmente
necesaria cuando la bibliografía se gestiona en bases sincronizadas que
pueden conservar entradas retractadas.

\subsection{Trabajo futuro}

Cuatro extensiones quedan planteadas sobre el marco construido: (i) cerrar
los balances de nitrógeno y azufre, que convertirían NH$_3$ y H$_2$S de
variables de supervisión en estados del modelo; (ii) reemplazar el supuesto
de población homogénea por una ecuación de estructura poblacional
(Observación~\ref{rem:homogeneidad}); (iii) la paralelización por
multiplexación de gases (Observación~\ref{rem:paralelizacion}), que
convertiría las réplicas secuenciales en simultáneas; y (iv) la dependencia
térmica explícita de los parámetros fisiológicos (factores $Q_{10}$ o
Arrhenius, supuesto B4), necesaria para operar el gemelo fuera de la zona
de confort. La simulación CFD de la línea neumática, anunciada en la
Sección~\ref{sec:diseno}, es el complemento experimental inmediato para
verificar el supuesto A1.


\bibliographystyle{plainnat}
\bibliography{refs}

\end{document}
